\section{Code Generation}

Code generation has long been a fundamental topic in computer science and software engineering, concerned with automating the process of producing code from specifications or descriptions of intent. Traditional approaches, including template-based methods, domain-specific languages, and rule-based systems, provided limited flexibility and adaptability to diverse programming tasks. With the rapid development of large language models in recent years, a new paradigm of neural code generation has emerged, leveraging the semantic understanding and linguistic capabilities of LLMs to automatically produce functional code from natural language descriptions. These models demonstrate strong code comprehension and writing abilities across multiple programming languages and paradigms, enabling application to a wide variety of software engineering tasks and significantly enhancing developer productivity \cite{10403378}.


\subsection{LLMs for Code Generation}

Large language models have rapidly advanced code generation through a multi-stage process. Prior to training, data preprocessing ensures that datasets from open-source repositories are clean, standardized, and suitable for model learning. During training, LLMs construct complex internal representations of syntax, semantics, and interrelations among code elements. Post-training, LLMs generate code by analyzing prompts, retrieving relevant patterns, assembling code fragments, and producing final outputs \cite{huynh2025largelanguagemodelscode}. Notable milestones including GitHub Copilot (2021) and Replit Ghostwriter (2022) demonstrated the increasing ability of LLMs to complete, explain, transform, and debug code.

As of 2025, leading models including OpenAI's Codex, Anthropic's Claude Sonnet 4.5, Google's Gemini 2.5 Pro, and Alibaba's Qwen3-Coder have demonstrated substantial improvements in reasoning, long-context understanding, and autonomous problem-solving capabilities \cite{openai2025codex, anthropic2025claude, google2025gemini, alibaba2025qwen}. These models excel in generating, debugging, and refactoring complex multi-file codebases with minimal human guidance. 

Despite these advances, significant limitations persist. LLMs require substantial computational resources, with large models demanding millions of GPU hours for training \cite{huynh2025largelanguagemodelscode}. They remain vulnerable to syntactic and semantic errors that can produce incorrect or non-functional code, and bias in generated code is a documented concern, with studies reporting gender-biased outputs in a notable proportion of cases. Security risks also arise from unsanitized training data, potentially introducing vulnerabilities such as prompt injection or memory safety issues. These challenges underscore the need for careful evaluation and mitigation strategies to ensure that LLM-based code generation is reliable, fair, and secure.
